\documentclass[11pt]{article}

\usepackage[margin=1in]{geometry}
\usepackage{amsmath,amssymb,amsthm}
\usepackage{graphicx}
\usepackage[hidelinks]{hyperref}

\newtheorem{theorem}{Theorem}
\newtheorem{lemma}{Lemma}
\newtheorem{remark}{Remark}

\newcommand{\norm}[1]{\left\lVert #1 \right\rVert}
\newcommand{\abs}[1]{\left\lvert #1 \right\rvert}
\newcommand{\Lzero}{L_0}
\newcommand{\FremlinTensor}{\widehat{\otimes}_{|\pi|}}

\title{A Strictly Positive Functional Subcase for a Kothe-Space Fremlin Tensor Representation Question}
\author{Automated Math Discovery Smoke Test}
\date{June 28, 2026}

\begin{document}
\maketitle

\section*{Source and Open Question}

This packet concerns the function-space representation remark in Section 4
of:

\begin{quote}
Qingying Bu, Gerard Buskes, Alexey I. Popov, Adi Tcaciuc, and
Vladimir G. Troitsky, \emph{The 2-concavification of a Banach lattice
equals the diagonal of the Fremlin tensor square}, arXiv:1109.5380.
\end{quote}

On page 13 the source paper states:

\begin{center}
\includegraphics[width=\textwidth]{figures/open_problem_crop.png}
\end{center}

In words, if \(E\) is a Kothe space on a sigma-finite measure space
\((\Omega,\Sigma,\mu)\), then \(E\otimes E\) is represented in
\(\Lzero(\Omega^2)\) by
\[
  (x\otimes y)(s,t)=x(s)y(t).
\]
The authors say that they do not know whether the completed Fremlin tensor
product \(E\FremlinTensor E\) can still be viewed as a sublattice of
\(\Lzero(\Omega^2)\), although this holds in many important special cases.

\section*{Candidate Partial Result}

\begin{theorem}
Let \(E\) be a Kothe Banach lattice on a sigma-finite measure space
\((\Omega,\Sigma,\mu)\). Suppose there is a measurable function
\(h>0\) a.e. and a constant \(C<\infty\) such that
\[
  \int_\Omega \abs{x(\omega)}h(\omega)\,d\mu(\omega)
  \le C\norm{x}_E
  \qquad (x\in E).
\]
Then the canonical representation
\[
  x\otimes y \longmapsto \bigl[(s,t)\mapsto x(s)y(t)\bigr]
\]
extends to an injective lattice homomorphism
\[
  J:E\FremlinTensor E\longrightarrow \Lzero(\Omega^2,\mu\otimes\mu).
\]
Consequently, under this strictly positive Kothe-dual functional
hypothesis, \(E\FremlinTensor E\) can be viewed as a vector sublattice of
\(\Lzero(\Omega^2)\) extending the usual elementary-tensor
representation.
\end{theorem}

\section*{Proof Intuition}

The obstacle in the source question is that the Fremlin norm completion
could, a priori, create abstract lattice elements with no measurable
representative on the original product space. A strictly positive
continuous Kothe-dual functional prevents this. It gives an equivalent
weighted product measure \(h(s)h(t)\,d\mu(s)d\mu(t)\) on \(\Omega^2\).
The elementary product map is then not merely a map into \(\Lzero\); it is
a bounded lattice bimorphism into the honest Banach lattice
\(L_1(h\otimes h)\). Fremlin's universal property extends this map to the
completed tensor product. Strict positivity plus Fremlin's positive-cone
domination then rules out a kernel.

\section*{Proof}

Define
\[
  d\nu(s,t)=h(s)h(t)\,d\mu(s)d\mu(t)
\]
on \(\Omega^2\). Since \(h>0\) a.e., the measures \(\nu\) and
\(\mu\otimes\mu\) have the same null sets.

Consider the bilinear map
\[
  m:E\times E\longrightarrow L_1(\nu),\qquad
  m(x,y)(s,t)=x(s)y(t).
\]
It is positive, and it is a lattice bimorphism because
\[
  \abs{m(x,y)}=m(\abs{x},\abs{y})
  \qquad (x,y\in E).
\]
Moreover,
\[
  \norm{m(x,y)}_{L_1(\nu)}
  =
  \int_{\Omega^2}\abs{x(s)y(t)}h(s)h(t)\,d\mu(s)d\mu(t)
  =
  \left(\int_\Omega\abs{x}h\,d\mu\right)
  \left(\int_\Omega\abs{y}h\,d\mu\right)
  \le C^2\norm{x}_E\norm{y}_E .
\]
Thus \(m\) is a continuous positive lattice bimorphism into the Banach
lattice \(L_1(\nu)\).

By the universal property of the Fremlin projective tensor product
\cite[1E(iii), 1F]{Fremlin74}, \(m\) induces a bounded positive operator
\[
  J:E\FremlinTensor E\longrightarrow L_1(\nu)
\]
such that \(J(x\otimes y)=m(x,y)\). Since \(m\) is a lattice bimorphism,
the same universal property gives that \(J\) is a lattice homomorphism.
Composing with the canonical inclusion \(L_1(\nu)\subset \Lzero(\Omega^2)\)
gives the desired extension into \(\Lzero(\Omega^2)\).

It remains only to prove injectivity. Suppose \(0<U\in E\FremlinTensor E\).
Fremlin's positive-cone domination property for the projective tensor
product says that there are nonzero \(x,y\in E_+\) with
\[
  0<x\otimes y\le U
\]
in \(E\FremlinTensor E\); see \cite[1(A)d]{Fremlin74}. Since \(x\) and
\(y\) are nonzero positive functions and \(h>0\) a.e., the product
\((s,t)\mapsto x(s)y(t)\) is nonzero in \(L_1(\nu)\). Therefore
\[
  0<J(x\otimes y)\le J(U),
\]
so \(J(U)\neq 0\). Hence no nonzero positive element lies in the kernel
of \(J\).

Finally, if \(J(W)=0\) for an arbitrary \(W\in E\FremlinTensor E\), then
\[
  J(\abs{W})=\abs{J(W)}=0.
\]
The previous paragraph applied to the positive element \(\abs{W}\) gives
\(\abs{W}=0\), and hence \(W=0\). Thus \(J\) is injective.

\section*{Verification Notes}

No computational verifier was needed. The proof is a direct universal
property argument. The most important checks for human review are:

\begin{itemize}
\item the weighted \(L_1\) estimate, which is exactly where the continuous
strictly positive Kothe-dual functional is used;
\item the fact that the pointwise product map is a lattice bimorphism;
\item the citation to Fremlin's positive-cone domination property, used to
prove that the induced lattice homomorphism has trivial kernel.
\end{itemize}

\section*{Limitations and Novelty Check}

This is a positive subcase, not a complete answer to the source question as
stated. The packet assumes the existence of a strictly positive continuous
Kothe-dual functional \(h\). This hypothesis holds in the standard examples
where such an associate weight is available, but the packet does not prove
that every Kothe space in the minimal source-paper definition has this
property.

The bounded literature check searched the run indexes for arXiv:1109.5380,
``2-concavification'', ``Fremlin tensor'', ``Kothe'', and \(L_0\) product
phrases. It also checked the later related papers arXiv:1301.0749,
arXiv:2404.02851, and arXiv:2601.19661. The 2013 Troitsky--Zabeti paper
extends the diagonal/concavification theorem, while the 2024 Zabeti paper
gives representation results in \(S(X\times Y)\)-type spaces for order
continuous Banach lattices. Neither source was found to give this exact
same-\(\Omega^2\), weighted-\(L_1\) sufficient condition as a resolution of
the source remark. Web searches on June 28, 2026 for exact phrases from the
remark and for Fremlin tensor products of Banach function/Kothe spaces did
not reveal an explicit full answer to the unrestricted question.

\section*{References}

\begin{thebibliography}{9}

\bibitem{BBPTT}
Q. Bu, G. Buskes, A. I. Popov, A. Tcaciuc, and V. G. Troitsky,
\emph{The 2-concavification of a Banach lattice equals the diagonal of
the Fremlin tensor square}, arXiv:1109.5380.

\bibitem{Fremlin74}
D. H. Fremlin,
\emph{Tensor products of Banach lattices},
Math. Ann. \textbf{211} (1974), 87--106.

\bibitem{TZ13}
V. G. Troitsky and O. Zabeti,
\emph{Fremlin tensor products of concavifications of Banach lattices},
arXiv:1301.0749.

\bibitem{Zabeti24}
O. Zabeti,
\emph{A representation theorem for order continuous Banach lattices with
applications}, arXiv:2404.02851.

\end{thebibliography}

\end{document}
